\documentclass[11pt]{article}
\usepackage{../latex/styles/core} 

\bibliography{bibliography}

\title{OSTEP Notes}
\author{Sai Mohan Kilambi}
\date{\today}
\begin{document}

\maketitle

\setcounter{tocdepth}{4}
\tableofcontents

\newpage
\section{Mechanism - Limited Direct Execution}

How can we virtualize the CPU efficiently?

\section{Book Chapter 14: Interlude - Memory API}
\subsection{Notes}

Main idea is there are two types of memory:
\begin{enumerate}
	\item Stack: Managed by compiler for you. Also called automatic memory. The compiler automatically allocates and deallocates. So if you want some data to persist then dont put it on stack.
	\item Heap: You the designer manages it. Use \textbf{malloc} to assign.
\end{enumerate}


\begin{lstlisting}[style=cppstyle]
	
// Chapter 14: Interlude - Memory API

#include <stdio.h>
#include <stdlib.h>

int main() {
  int *a;

  // This line has both allocation on stack and heap
  // The pointer 'a' is allocated on the stack
  // The memory for 10 integers is allocated on the heap
  a = (int *)malloc(10 * sizeof(int));

  // Print the address of the allocated memory
  // This will print the address of where the 10 integers are stored on the heap
  // Its also the first element of the array when you dereference it
  printf("%p\n", a);

  // Free the allocated memory
  free(a);
  return 0;
}
\end{lstlisting}
\subsection{Chapter End Exercises}
\subsubsection{Null.c}

We basically created a null.c program and tried to dereference a pointer that has address as $NULL$. We saw a segfault. Makes sense we are trying to access a memory location we are not allowed. 

\begin{lstlisting}[style=cppstyle]
#include <iostream>
#include <unistd.h>
using namespace std;
int main() {
  int *a;
  a = NULL;
  cout << "a: " << a << " *a: " << *a << endl;
  return 0;
}
\end{lstlisting}

\subsubsection{Add -g compiler flag}
The $-g$. This create the $dSYM$ directory with symbols annotated. This way $lldb$ can used to debug as it has the all the info. To debug you can do
\begin{lstlisting}[style=asmstyle]

	lldb null
	
\end{lstlisting}
This will open the debugger. \lstinline[style=bashstyle]|run| can be used to run the code. It will crash at the deferenccing line.

\section{XV6}

\subsection{Operating System Interfaces}
OS provides services to user programs through interfaces. We want the interfaces to be simple on one hand but also provide generality and there is always tension in this. 

User programs execute as \textbf{processes}. When the process needs to run a privileged instruction, it makes a \textbf{system call}. This moves it into kernel space. The kernel executes the requested service and then returns back into the user space. Some important system calls 
\begin{enumerate}
	\item fork
	\item exec 
\end{enumerate}

If you look at $/usr/sh.c$ you will find that its really a very simple program. When you open a terminal thats a GUI and the first thing it starts doing is running the shell executable. This executable is just waiting for the user to type in a command. The shell program forks and if the fork is successful will start executing this new process. 

XV6 uses the \textbf{ELF} format to execute instructions. It has a certain format for instructions and data. 

How do you think fork and exec work in conjunction? Think

\subsection{Traps and System Calls (Chapter 4 of XV6 book)}

Three events that cause the CPU to set aside ordinary execution and force transfer to special code that handles the event.
\begin{enumerate}
	\item System Calls: \textbf{ecall} instruction
	\item Exceptions: Some illegal thing was attempted like divide by 0.
	\item Interrupt
\end{enumerate}

We use the term \textbf{trap} generically for this handling of events.
It proceeds in four stages.
\begin{enumerate}
	\item Hardware actions taken by the RISC-V CPU.
	\item Some assembly instructions that prepare way for the Kernel C code.
	\item A C function that decides what to do with the trap.
	\item The system call or the device driver service routine.
\end{enumerate}


\section{xv6}

\subsection{Remzi's Video \cite{xv6_remzi}}
The main idea is how to go from \textbf{user} mode to \textbf{kernel} mode. Every xv6 user program exits with a \textbf{exit()} command. 

We added a hello.c command in $user/$. I compiled it with the command.

\begin{lstlisting}[style=bashstyle]
	make qemu #OR
	make fs.img
\end{lstlisting}


This will log into xv6. To exit the key stroke is \textbf{C-a a x}

Note when you compiled three things happened
\begin{lstlisting}[style=bashstyle]
	riscv64-unknown-elf-gcc -Wall -Werror -Wno-unknown-attributes -O -fno-omit-frame-pointer -ggdb -gdwarf-2 -march=rv64gc -MD -mcmodel=medany -ffreestanding -fno-common -nostdlib -fno-builtin-strncpy -fno-builtin-strncmp -fno-builtin-strlen -fno-builtin-memset -fno-builtin-memmove -fno-builtin-memcmp -fno-builtin-log -fno-builtin-bzero -fno-builtin-strchr -fno-builtin-exit -fno-builtin-malloc -fno-builtin-putc -fno-builtin-free -fno-builtin-memcpy -Wno-main -fno-builtin-printf -fno-builtin-fprintf -fno-builtin-vprintf -I. -fno-stack-protector -fno-pie -no-pie   -c -o user/hello.o user/hello.c
	
	riscv64-unknown-elf-ld -z max-page-size=4096 -T user/user.ld -o user/_hello user/hello.o user/ulib.o user/usys.o user/printf.o user/umalloc.o

	riscv64-unknown-elf-objdump -S user/_hello > user/hello.asm

	riscv64-unknown-elf-objdump -t user/_hello | sed '1,/SYMBOL TABLE/d; s/ .* / /; /^$/d' > user/hello.sym
\end{lstlisting}

The ld portion is called the linker. And its linking your code to other object files like \textbf{printf}.

If you take a look through the \textbf{usys.S} file, its assembly instructions. Heres a chunk:
\begin{lstlisting}[style=asmstyle]
	.global uptime // Make it available for the linker to see.
	uptime: 
 	li a7, SYS_uptime // a7 is the register that gets the syscall number
 	ecall // Trigger supervisor kernel. When the kernel runs it looks at the a7 register and jumps to this call.
 	ret // Return from the kernel to the caller
\end{lstlisting}

One thing you might wonder, where are the arguments? Like $open$ will have arguments. The $xv6$ convention is to put them on the stack for the kernel to pop off. 

\printbibliography

\end{document}
